\section{Full-Order Numerical Model}

The governing energy equation is discretized in space to obtain
a system of ordinary differential equations, which is integrated
iteratively to reach the steady-state solution. This section
describes the computational grid, spatial discretization,
boundary implementation, and numerical solution procedure.

\subsection{Computational Grid and Degrees of Freedom}

A uniform Cartesian grid is used to discretize the rectangular
computational domain. The grid consists of

\begin{equation}
    N_x=50,
    \qquad
    N_y=20,
\end{equation}

giving a total of

\begin{equation}
    \boxed{N=N_xN_y=1000}
\end{equation}

physical degrees of freedom.

The grid spacings in the streamwise and wall-normal directions
are

\begin{equation}
    \Delta x=\frac{L_x}{N_x}=0.1,
    \qquad
    \Delta y=\frac{L_y}{N_y}=0.05.
\end{equation}

A cell-centred arrangement is employed, with the physical cell
centres located at

\begin{equation}
    x_i=\left(i-\frac{1}{2}\right)\Delta x,
    \qquad
    y_j=\left(j-\frac{1}{2}\right)\Delta y,
\end{equation}

where $i=1,\ldots,N_x$ and $j=1,\ldots,N_y$.

The physical temperature field is represented by a two-dimensional
array of size $N_x\times N_y$, or equivalently by a vector
containing $N=1000$ unknowns. One layer of ghost cells is
allocated around the physical domain to implement the boundary
conditions. Therefore, the extended temperature array has
dimensions

\begin{equation}
    (N_x+2)\times(N_y+2)=52\times22.
\end{equation}

For the algebraic formulation, the physical temperature field is
converted into a one-dimensional vector using C-order
vectorization, in which the wall-normal index varies fastest.
The mapping between the two-dimensional indices and the vector
index is

\begin{equation}
    k=(i-1)N_y+(j-1).
    \label{eq:vector_indexing}
\end{equation}

Thus, the full-order temperature vector is written as

\begin{equation}
    \mathbf{T}\in\mathbb{R}^{1000}.
\end{equation}

The ghost-cell values are used in the spatial discretization
but are not included among the physical degrees of freedom.


\subsection{Spatial Discretization}

The spatial derivatives are approximated using second-order
central differences on the uniform Cartesian grid. The
discretized energy equation can be expressed as a system of
ordinary differential equations:

\begin{equation}
    \frac{d\mathbf{T}}{dt}
    =
    \mathcal{L}(\mathbf{T}),
    \label{eq:direct_ode}
\end{equation}

where $\mathbf{T}$ is the vector of cell-centred temperatures
and $\mathcal{L}$ denotes the spatial operator, including
convection, diffusion, viscous dissipation, and boundary
contributions.

For the direct solver, the spatial operator at cell $(i,j)$
is expressed in conservative form as

\begin{align}
    \mathcal{L}_{ij}(\mathbf{T})
    ={}&
    -c\frac{
    u_{i+1,j}T_{i+1,j}
    -u_{i-1,j}T_{i-1,j}
    }{2\Delta x}
    \nonumber\\
    &-
    c\frac{
    v_{i,j+1}T_{i,j+1}
    -v_{i,j-1}T_{i,j-1}
    }{2\Delta y}
    \nonumber\\
    &+
    \frac{1}{Pe}
    \frac{
    T_{i+1,j}-2T_{i,j}+T_{i-1,j}
    }{\Delta x^2}
    \nonumber\\
    &+
    \frac{1}{Pe}
    \frac{
    T_{i,j+1}-2T_{i,j}+T_{i,j-1}
    }{\Delta y^2}
    +S_{ij}.
    \label{eq:direct_spatial_operator}
\end{align}

The first two terms represent convective transport in the
streamwise and wall-normal directions, respectively. The
next two terms approximate thermal diffusion in the
corresponding directions. All spatial derivatives are
discretized using second-order central differences.

The source term $S_{ij}$ accounts for viscous dissipation.
It is evaluated from the prescribed velocity field and its
discrete spatial derivatives:

\begin{equation}
\begin{aligned}
    S_{ij}
    =\frac{\mathrm{Ec}}{\mathrm{Re}}
    \Bigg[
    &2\left(
    \frac{u_{i+1,j}-u_{i-1,j}}{2\Delta x}
    \right)^2
    \\[4pt]
    &+2\left(
    \frac{v_{i,j+1}-v_{i,j-1}}{2\Delta y}
    \right)^2
    \\[4pt]
    &+\left(
    \frac{u_{i,j+1}-u_{i,j-1}}{2\Delta y}
    +
    \frac{v_{i+1,j}-v_{i-1,j}}{2\Delta x}
    \right)^2
    \Bigg].
\end{aligned}
\label{eq:discrete_viscous_source}
\end{equation}

Equivalently, using the parameter definitions introduced in
Section~2.2, the source term can be expressed as

\begin{equation}
    S_{ij}=\frac{Br}{Pe}\Phi_{ij}.
\end{equation}

Since the velocity field and viscous-dissipation source are
prescribed and independent of temperature, the spatial operator
is linear in the temperature unknowns, provided that the
boundary conditions are linear.

After incorporating the boundary conditions, the resulting
discrete system takes the affine form

\begin{equation}
    \boxed{
    \frac{d\mathbf{T}}{dt}
    =\mathbf{A}\mathbf{T}+\mathbf{b},
    }
    \label{eq:semi_discrete_system}
\end{equation}

where $\mathbf{A}\in\mathbb{R}^{N\times N}$ is the sparse
spatial matrix and $\mathbf{b}\in\mathbb{R}^{N}$ is the
affine forcing vector.

The matrix $\mathbf{A}$ contains the discretized convective
and diffusive contributions. The vector $\mathbf{b}$ accounts
for the prescribed viscous-dissipation source and constant
contributions arising from the nonhomogeneous boundary
conditions.

The spatial matrix is assembled using a five-point stencil,
in which each physical cell interacts with itself and its
four immediate neighbours. After vectorization, the resulting
matrix has a sparse structure. Its coefficients are modified
near the boundaries to account for the elimination of ghost
cells.

For the prescribed velocity field, $u=u(y)$ and $v=0$, the
conservative convective discretization is consistent with
the non-conservative form used in the steady-state Galerkin
formulation, since the discrete streamwise velocity derivative
and wall-normal velocity derivative vanish.


\subsection{Ghost-Cell Boundary Implementation}

The boundary conditions introduced in Section~2.3 are
implemented through one layer of ghost cells surrounding
the physical domain. This approach allows the same central-difference
stencils to be applied to physical cells adjacent to the
boundaries.

For a Dirichlet boundary condition, the ghost-cell temperature
is obtained by reflecting the adjacent physical-cell temperature
about the prescribed boundary temperature:

\begin{equation}
    \boxed{
    T_g=2T_w-T_{\mathrm{adjacent}},
    }
    \label{eq:dirichlet_ghost}
\end{equation}

where $T_w$ is the prescribed wall or boundary temperature.

For the inlet, bottom wall, and top wall, the corresponding
relations are

\begin{align}
    T_{0,j}&=2y_j-T_{1,j},\\
    T_{i,0}&=-T_{i,1},\\
    T_{i,N_y+1}&=2-T_{i,N_y}.
\end{align}

At the outlet, the zero streamwise temperature gradient is
implemented using a zero-gradient ghost-cell relation:

\begin{equation}
    \boxed{
    T_{N_x+1,j}=T_{N_x,j}.
    }
    \label{eq:neumann_ghost}
\end{equation}

The ghost-cell values are substituted into the spatial
finite-difference expressions and eliminated from the
algebraic system. For Dirichlet boundaries, this substitution
modifies the coefficients associated with the adjacent
physical cells and introduces constant contributions into
the forcing vector. The zero-gradient outlet condition
modifies the corresponding boundary-row coefficients
without introducing a prescribed-temperature contribution.

Consequently, the boundary conditions are incorporated
directly into the assembled matrix and forcing vector:

\begin{equation}
    \frac{d\mathbf{T}}{dt}
    =
    \mathbf{A}\mathbf{T}+\mathbf{b}.
\end{equation}

The same discrete boundary treatment is retained in the
Galerkin reduced-order formulation, ensuring consistency
between the full-order and reduced-order models.

% Suggested figure:
% Physical cells and ghost cells at the four boundaries.
%
% \begin{figure}[htbp]
%     \centering
%     \includegraphics[width=0.75\textwidth]
%     {figures/ghost_cell_boundary_conditions.pdf}
%     \caption{Ghost-cell arrangement and boundary treatment
%     for the full-order numerical model.}
%     \label{fig:ghost_cells}
% \end{figure}


\subsection{Steady-State Solution Procedure}

The steady-state solution is obtained by integrating the
semi-discrete system in pseudo-time using the fully implicit
backward-Euler method. Applying this time discretization to
Eq.~\eqref{eq:semi_discrete_system} gives

\begin{equation}
    \frac{
    \mathbf{T}^{n+1}-\mathbf{T}^{n}
    }{\Delta t}
    =
    \mathbf{A}\mathbf{T}^{n+1}+\mathbf{b}.
\end{equation}

Rearranging, the temperature field at the next iteration
is obtained from

\begin{equation}
    \boxed{
    \left(
    \mathbf{I}-\Delta t\mathbf{A}
    \right)\mathbf{T}^{n+1}
    =
    \mathbf{T}^{n}+\Delta t\mathbf{b}.
    }
    \label{eq:backward_euler}
\end{equation}

The numerical solution procedure consists of the following
steps:

\begin{enumerate}
    \item Assemble the sparse spatial matrix and affine
    forcing vector for the specified parameter combination.

    \item Construct the implicit system matrix
    $\mathbf{I}-\Delta t\mathbf{A}$.

    \item Perform sparse LU factorization of the system matrix.

    \item Reuse the LU factorization during subsequent
    iterations for the same parameter combination.

    \item Solve for the updated physical temperature field.

    \item Update the ghost-cell values according to the
    prescribed boundary conditions.

    \item Evaluate the steady-state residual.

    \item Repeat the procedure until the convergence
    criterion is satisfied or the maximum iteration count
    is reached.
\end{enumerate}

Since the spatial matrix and forcing vector remain constant
during the iteration for a fixed parameter combination,
the LU factorization can be reused. A new factorization
is required when the governing parameters change and
the system matrix is reconstructed.

It should be noted that the time-marching procedure is
used only as a numerical means of obtaining the steady-state
solution. The reduced-order model developed in this work
is based on the steady-state governing equation rather
than the transient time-dependent system.


\subsection{Convergence Criterion}

Convergence is assessed using the root-mean-square (RMS)
of the discrete steady-state residual. For the semi-discrete
system, the residual at iteration $n$ is defined as

\begin{equation}
    \mathbf{R}^{n}
    =
    \mathbf{A}\mathbf{T}^{n}+\mathbf{b}.
\end{equation}

The corresponding RMS residual is

\begin{equation}
    \boxed{
    R_{\mathrm{RMS}}^{n}
    =
    \sqrt{
    \frac{1}{N}
    \sum_{i=1}^{N}
    \left(R_i^n\right)^2
    }.
    }
    \label{eq:fom_residual_rms}
\end{equation}

The iteration is terminated when the residual satisfies

\begin{equation}
    R_{\mathrm{RMS}}^{n}
    \leq \varepsilon_{\mathrm{FOM}},
\end{equation}

where $\varepsilon_{\mathrm{FOM}}$ is the prescribed
convergence tolerance. The maximum number of iterations
is imposed as an additional stopping condition.

The steady-state residual measures how closely the
numerical solution satisfies the discretized governing
equation. It is distinct from the reduced-order model
approximation error, which measures the difference
between the full-order and reconstructed reduced-order
temperature fields and is discussed in Section~6.


\subsection{Numerical Settings and Implementation}

The principal numerical settings of the full-order solver
are summarized in Table~\ref{tab:fom_settings}.

\begin{table}[htbp]
    \centering
    \caption{Numerical settings of the full-order model.}
    \label{tab:fom_settings}
    \begin{tabular}{lll}
        \toprule
        \textbf{Parameter} &
        \textbf{Description} &
        \textbf{Value} \\
        \midrule
        $N_x$ & Number of streamwise cells & 50 \\
        $N_y$ & Number of wall-normal cells & 20 \\
        $N$ & Number of physical unknowns & 1000 \\
        $\Delta x$ & Streamwise grid spacing & 0.1 \\
        $\Delta y$ & Wall-normal grid spacing & 0.05 \\
        $\Delta t$ & Pseudo-time step & 1 \\
        $\varepsilon_{\mathrm{FOM}}$ & RMS residual tolerance & $10^{-10}$ \\
        $N_{\max}$ & Maximum iterations & $10^5$ \\
        \bottomrule
    \end{tabular}
\end{table}

The implementation of the full-order numerical solver,
including the spatial discretization, boundary treatment,
and solution procedure, is available in the project
repository:

\begin{center}
    \href{
    https://github.com/Arian-Shahnoori/CFD/tree/main/CFD_Project_03
    }{
    Full-order numerical solver -- GitHub repository
    }
\end{center}

\subsection{Code implementation}

The full-order solver is implemented in Python, with its main
numerical routines contained in \texttt{src/solvers/direct.py}.
The implementation is organized into four functions:

\begin{itemize}
    \item \texttt{build\_spatial\_matrix}: Constructs the sparse
    five-diagonal matrix $\mathbf{A}$. It maps the two-dimensional
    physical cell indices to the one-dimensional system indices,
    calculates the coefficients of the five-point stencil, and
    incorporates the ghost-cell boundary conditions into the
    corresponding matrix coefficients.

    \item \texttt{build\_rhs}: Constructs the right-hand-side
    contribution for a single row. This function was originally
    developed for the Alternating Plane Filling (APF) method.

    \item \texttt{build\_forcing\_field}: Applies
    \texttt{build\_rhs} across the computational domain to
    assemble the complete forcing vector $\mathbf{b}$.

    \item \texttt{solve\_with\_direct}: Assembles the spatial
    matrix and forcing vector, factorizes
    $(\mathbf{I}-\Delta t\mathbf{A})$ using sparse LU
    factorization, and advances the temperature field.
    The steady-state residual is evaluated using
    \texttt{steady\_state\_residual\_rms} to determine
    convergence.
\end{itemize}

This modular structure separates spatial discretization,
forcing-term construction, and numerical solution. The same
discrete spatial formulation is subsequently used to construct
the Galerkin reduced-order model.